\documentclass[11pt,a4paper]{article}

\usepackage[utf8]{inputenc}
\usepackage[T1]{fontenc}
\usepackage[spanish,es-tabla,es-noshorthands]{babel}
\usepackage{amsmath,amssymb,amsthm}
\usepackage[margin=2.5cm]{geometry}
\usepackage{enumitem}
\usepackage{hyperref}
\usepackage{xcolor}
\usepackage{tcolorbox}
\usepackage{array}
\usepackage{booktabs}
\usepackage{parskip}

\hypersetup{
    colorlinks=true,
    linkcolor=blue!60!black,
    urlcolor=blue!60!black,
    citecolor=blue!60!black
}

\newtheoremstyle{definicion}
    {1em}{1em}{}{}{\bfseries}{.}{.5em}{}
\theoremstyle{definicion}
\newtheorem{definicion}{Definición}[section]
\newtheorem{teorema}[definicion]{Teorema}
\newtheorem{proposicion}[definicion]{Proposición}
\newtheorem{ejemplo}[definicion]{Ejemplo}
\newtheorem{observacion}[definicion]{Observación}

\newtcolorbox{intuicion}{
    colback=blue!5!white,
    colframe=blue!60!black,
    title=Intuición,
    fonttitle=\bfseries,
    boxrule=0.5pt,
    arc=2pt
}

\newtcolorbox{idea}{
    colback=green!5!white,
    colframe=green!40!black,
    title=Idea clave,
    fonttitle=\bfseries,
    boxrule=0.5pt,
    arc=2pt
}

\newcommand{\MT}{\textsc{mt}}
\newcommand{\MTs}{\textsc{mt}s}
\newcommand{\enc}[1]{\langle #1 \rangle}
\newcommand{\Lang}{\mathcal{L}}
\newcommand{\N}{\mathbb{N}}
\newcommand{\R}{\mathbb{R}}

\title{\textbf{Resumen Teórico: Complejidad Temporal}\\[0.3em]
\large Computabilidad y Complejidad --- UNRC}
\author{Basado en el material de Pablo Castro (UNRC)}
\date{}

\begin{document}

\maketitle

\tableofcontents

\vspace{1em}
\hrule
\vspace{1em}

\section{Introducción}

Hasta ahora vimos qué problemas se pueden \emph{computar} y cuáles no. El estudio de la \textbf{computabilidad} responde a la pregunta:
\begin{quote}
\emph{Dado un problema $P$, ¿es computable?}
\end{quote}

La \textbf{complejidad computacional} cambia el enfoque: ya trabajamos solo con problemas \emph{computables} y nos preguntamos:
\begin{quote}
\emph{¿Hay algoritmos \textbf{eficientes} para resolver $P$?}
\end{quote}

Este resumen recorre las herramientas básicas para responder esta pregunta:
\begin{enumerate}
    \item Cómo medir el tiempo de ejecución de un algoritmo.
    \item Las notaciones asintóticas $O, o, \Omega, \Theta$.
    \item Las clases de complejidad $\mathrm{Time}(t(n))$ y, en particular, la clase $P$.
    \item Cómo se relacionan los distintos modelos de computación bajo la lupa temporal: \MTs\ con una cinta, multicinta y no determinista.
    \item La \emph{Tesis de Church-Turing extendida}: todos los modelos razonables son equivalentes en tiempo \emph{módulo polinomios}.
    \item Ejemplos clásicos de problemas en $P$: caminos en grafos, números coprimos y pertenencia a gramáticas libres de contexto.
\end{enumerate}

\begin{intuicion}
Una pregunta importante antes de hablar de eficiencia: \emph{¿qué modelo de computación usamos como referencia?} La respuesta es: las máquinas de Turing deterministas. Veremos que los modelos razonables se simulan entre sí con un costo \emph{a lo sumo polinomial}, así que la noción de ``algoritmo eficiente'' es robusta.
\end{intuicion}

\section{Tiempo de ejecución como función}

Dada una máquina de Turing $M$ que termina en todas sus entradas, su \textbf{tiempo de ejecución} es una función
\[
T_M : \N \to \R^{\geq 0},
\]
donde $T_M(n)$ representa la cantidad máxima de pasos que realiza $M$ para una entrada de tamaño $n$ (análisis del \textbf{peor caso}).

\textbf{Convenciones habituales:}
\begin{itemize}
    \item $T_M$ es \textbf{monótona} (no decrece con $n$).
    \item Trabajamos con \MTs\ que terminan siempre.
    \item Se analiza típicamente el \textbf{peor caso}, no el caso promedio.
    \item Nos interesa el \emph{comportamiento asintótico}: cómo crece $T_M$ cuando $n$ se hace grande.
\end{itemize}

\begin{intuicion}
No interesa el valor exacto de $T_M(7)$, sino cómo escala $T_M(n)$ cuando $n \to \infty$. Las constantes multiplicativas y los términos de menor orden no son lo importante: lo que importa es el \emph{orden de crecimiento}.
\end{intuicion}

\section{Notaciones asintóticas}

Para hablar del orden de crecimiento de las funciones se usan varias notaciones, todas ellas basadas en comparaciones asintóticas.

\subsection{Notación $O$ (big-oh): cota superior}

\begin{definicion}[Big-oh]
$f(n) \in O(g(n))$ si y solo si existen constantes $c > 0$ y $n_0 > 0$ tales que
\[
\forall n \geq n_0 : \quad f(n) \leq c \cdot g(n).
\]
\end{definicion}

\textbf{Intuición:} para entradas suficientemente grandes, $g$ crece \emph{al menos tan rápido} como $f$, salvo una constante multiplicativa.

\begin{ejemplo}
$3n^2 + 5n + 100 \in O(n^2)$. Tomando $c = 4$ y $n_0$ adecuado, $3n^2 + 5n + 100 \leq 4n^2$ a partir de cierto punto.
\end{ejemplo}

\subsection{Notación $o$ (little-oh): cota superior estricta}

\begin{definicion}[Little-oh]
$f(n) \in o(g(n))$ si y solo si para \emph{toda} constante $c > 0$ existe $n_0 > 0$ tal que
\[
\forall n \geq n_0 : \quad f(n) < c \cdot g(n).
\]
\end{definicion}

Una caracterización equivalente, muy útil, es mediante un límite:
\[
\lim_{n \to \infty} \frac{f(n)}{g(n)} = 0 \iff f(n) \in o(g(n)).
\]

\begin{intuicion}
$o$ es una versión \emph{estricta} de $O$: la función $g$ crece \emph{estrictamente} más rápido que $f$. En particular, $f(n) \notin o(f(n))$ (uno no crece estrictamente más rápido que sí mismo), aunque sí $f(n) \in O(f(n))$.
\end{intuicion}

\begin{ejemplo}
$n \in o(n^2)$, $\log n \in o(n)$, $n \in o(n \log n)$.
\end{ejemplo}

\subsection{Notación $\Omega$: cota inferior}

\begin{definicion}[Omega]
$f(n) \in \Omega(g(n))$ si y solo si existen constantes $c > 0$ y $n_0 > 0$ tales que
\[
\forall n \geq n_0 : \quad f(n) \geq c \cdot g(n).
\]
\end{definicion}

\textbf{Intuición:} para entradas grandes, $f$ crece \emph{al menos tan rápido} como $g$ (es la cota \emph{inferior}, opuesta a $O$).

\subsection{Notación $\Theta$: igualdad de orden}

\begin{definicion}[Theta]
$f(n) \in \Theta(g(n))$ si y solo si $f(n) \in O(g(n))$ \emph{y} $f(n) \in \Omega(g(n))$.
\end{definicion}

\textbf{Intuición:} $f$ y $g$ crecen ``al mismo ritmo'' --- una está sándwich entre dos múltiplos constantes de la otra.

\begin{ejemplo}
$3n^2 + 5n + 100 \in \Theta(n^2)$.
\end{ejemplo}

\begin{idea}
Resumen de las cuatro notaciones:
\begin{itemize}
    \item $O(g)$: ``$f$ no crece más rápido que $g$''.
    \item $o(g)$: ``$f$ crece \emph{estrictamente} más lento que $g$''.
    \item $\Omega(g)$: ``$f$ no crece más lento que $g$''.
    \item $\Theta(g)$: ``$f$ crece exactamente igual que $g$''.
\end{itemize}
\end{idea}

\section{Análisis de tiempo de ejecución: un ejemplo}

Estudiemos cuánto tarda una \MT\ en decidir el lenguaje clásico
\[
A = \{ 0^k 1^k \mid k \geq 0 \}.
\]

\subsection{Primera solución: $O(n^2)$}

\textbf{Algoritmo (alto nivel).} Dado el input $w$:
\begin{enumerate}
    \item Recorrer la cinta; si se encuentra ``10'', rechazar.
    \item Volver al principio. Tachar el primer 0, buscar el primer 1 y tacharlo.
    \item Repetir el paso 2 hasta que:
    \begin{itemize}
        \item No queden 0s ni 1s: \emph{aceptar}.
        \item Queden solo 0s o solo 1s: \emph{rechazar}.
    \end{itemize}
\end{enumerate}

\textbf{Análisis paso a paso} para una entrada de longitud $n$:
\begin{itemize}
    \item El paso 1 (verificar formato $0^*1^*$): $O(n)$ pasos.
    \item El paso 2 (tachar un par 0/1): cada iteración requiere recorrer la cinta, es decir $O(n)$ pasos.
    \item El paso 2 se repite hasta $n/2$ veces, una por cada par a tachar.
\end{itemize}

Tiempo total:
\[
O(n) + O\!\left(\frac{n}{2} \cdot n\right) = O(n) + O(n^2) = O(n^2).
\]

\textbf{Conclusión:} esta \MT\ decide $A$ en tiempo $O(n^2)$.

\subsection{Una solución mejor: $O(n \log n)$}

\textbf{Idea:} en cada pasada, eliminar la \emph{mitad} de los 0s y la mitad de los 1s. Si en alguna pasada la cantidad de 0s (o 1s) es impar y mayor a 1, rechazar.

\textbf{Algoritmo.}
\begin{enumerate}
    \item Recorrer la cinta; si aparece ``10'', rechazar.
    \item Repetir mientras queden 0s o 1s:
    \begin{itemize}
        \item Contar la paridad de 0s y 1s. Si alguno es impar (y mayor que 1), rechazar.
        \item Tachar uno de cada dos 0s alternadamente, y lo mismo con los 1s.
    \end{itemize}
    \item Si al terminar no quedan 0s ni 1s, aceptar; si no, rechazar.
\end{enumerate}

\textbf{Análisis.}
\begin{itemize}
    \item El paso 1: $O(n)$.
    \item El bucle del paso 2 se ejecuta $O(\log n)$ veces, porque en cada iteración \emph{se divide la cantidad por 2}.
    \item Cada iteración del bucle realiza $O(n)$ pasos (contar paridad y tachar alternadamente).
\end{itemize}

Tiempo total:
\[
O(n) + O(\log n) \cdot O(n) = O(n \log n).
\]

\begin{intuicion}
La mejora ($O(n^2) \to O(n \log n)$) se obtiene aprovechando que cada pasada \emph{divide a la mitad} el trabajo restante, igual que en \emph{binary search}. La idea es: ``cada vez hago la mitad'' produce un factor $\log n$ donde antes había un $n$.
\end{intuicion}

\begin{ejemplo}[Trazado para $n = 50$, $w = 0^{25} 1^{25}$]
\begin{itemize}
    \item Pasada 1: paridad OK; tacho 12 ceros y 12 unos. Quedan 13 ceros y 13 unos. Total tachado: 24+24.
    \item Pasada 2: paridad OK; tacho 6 ceros y 6 unos.
    \item Pasada 3: tacho 3 y 3.
    \item Pasada 4: tacho 2 y 2.
    \item Pasada 5: tacho 1 y 1. Quedan 0s. Aceptar.
\end{itemize}
Cantidad de pasadas: $\Theta(\log n)$.
\end{ejemplo}

\section{Clases de complejidad}

Habiendo definido cómo medir el tiempo, podemos \emph{clasificar} los problemas según el orden de crecimiento del algoritmo más rápido que los resuelve.

\begin{definicion}[Clase $\mathrm{Time}$]
\[
\mathrm{Time}(t(n)) = \{ L \mid L \text{ se puede decidir con una \MT\ en } O(t(n)) \text{ pasos}\}.
\]
\end{definicion}

\begin{ejemplo}
$\mathrm{Time}(O(n^2))$ contiene a todos los problemas que se pueden resolver con un algoritmo que en el peor caso hace $c \cdot n^2$ pasos para alguna constante $c$. Sorting es un ejemplo clásico.
\end{ejemplo}

\textbf{Preguntas centrales que estudia la complejidad:}
\begin{itemize}
    \item Dado un problema, ¿pertenece a $\mathrm{Time}(O(t(n)))$ para algún $t$ específico?
    \item ¿Hay problemas en $\mathrm{Time}(O(t(n)))$ pero no en $\mathrm{Time}(O(t'(n)))$ para $t' < t$?
\end{itemize}

\section{La clase $P$: problemas tratables}

Los algoritmos que corren en tiempo \emph{polinomial} reciben un nombre especial: se dicen \textbf{tratables}. La razón es eminentemente práctica.

\subsection{Comparación de órdenes de crecimiento}

La siguiente tabla, asumiendo $10^6$ operaciones por segundo, muestra cuán dramática es la diferencia entre los órdenes polinomial y exponencial:

\begin{center}
\begin{tabular}{lccccc}
\toprule
 & $n=10$ & $n=20$ & $n=30$ & $n=40$ & $n=50$ \\
\midrule
$n$    & $0.00001$ s & $0.00002$ s & $0.00003$ s & $0.00004$ s & $0.00005$ s \\
$n^2$  & $0.0002$ s  & $0.0004$ s  & $0.0009$ s  & $0.0016$ s  & $0.0025$ s \\
$n^3$  & $0.1$ s     & $3.2$ s     & $24.3$ s    & $1.7$ min   & $15.2$ min \\
$2^n$  & $0.001$ s   & $1.0$ s     & $17.9$ min  & $12.7$ días & $35.6$ años \\
\bottomrule
\end{tabular}
\end{center}

\begin{intuicion}
Los algoritmos polinomiales escalan bien: pasar de $n = 10$ a $n = 50$ multiplica el tiempo por algo manejable. Los algoritmos exponenciales colapsan: un problema que se resuelve en un segundo para $n = 20$ tarda \emph{décadas} para $n = 50$. Por eso identificamos ``tratable'' con ``polinomial'' a nivel teórico.
\end{intuicion}

\subsection{Definición de la clase $P$}

\begin{definicion}[Clase $P$]
\[
P = \bigcup_{k > 0} \mathrm{Time}(O(n^k)).
\]
Es decir, $P$ es la clase de los problemas para los cuales \emph{existe algún $k$} y un algoritmo polinomial de grado $k$ que los decide.
\end{definicion}

\textbf{Propiedades:}
\begin{itemize}
    \item $P$ es \textbf{invariante} respecto a los modelos de computación deterministas razonables (Tesis de Church-Turing Extendida).
    \item Se considera la clase de los problemas \emph{eficientes}: un avance en el hardware (multiplicar la velocidad por un factor constante) tiene impacto real sobre problemas polinomiales, pero apenas mueve la aguja para problemas exponenciales.
\end{itemize}

\section{Modelos de computación y simulación polinomial}

Una pregunta importante: ¿depende la clase $P$ del modelo de computación elegido?

Hay varios modelos relevantes:
\begin{itemize}
    \item \MTs\ deterministas con una sola cinta.
    \item \MTs\ deterministas con varias cintas.
    \item \MTs\ deterministas con cinta doblemente infinita.
    \item \MTs\ no deterministas.
\end{itemize}

\textbf{Resultado clave:} los \emph{tres primeros} (todos los modelos deterministas razonables) se simulan entre sí en \emph{tiempo polinomial}. Por lo tanto, definen la \emph{misma} clase $P$.

\subsection{Simulación: multicinta $\Rightarrow$ una sola cinta}

\begin{teorema}\label{teo:multi}
Si una \MT\ multicinta $M$ decide un problema en $O(t(n))$ pasos, entonces existe una \MT\ con una sola cinta $S$ que lo decide en $O(t(n)^2)$ pasos.
\end{teorema}

\textbf{Idea de la simulación.}
\begin{itemize}
    \item En la cinta única de $S$ se concatenan las $k$ cintas de $M$, separadas por un símbolo especial \texttt{\#}.
    \item Para cada cinta, se usa una marca sobre un símbolo (por ejemplo, $\overline{a}$) para indicar la posición del cabezal.
    \item Para simular \emph{un paso} de $M$: $S$ recorre todas sus celdas para localizar las $k$ cabezas marcadas, aplica la transición y vuelve a recorrer para actualizar los símbolos y mover los cabezales virtuales.
\end{itemize}

\textbf{Análisis del costo.}
\begin{itemize}
    \item Por cada paso de $M$, los contenidos virtuales tienen a lo sumo longitud total $O(t(n))$ (porque cada cinta no puede crecer más rápido que el tiempo gastado).
    \item Simular un paso de $M$ cuesta $O(t(n))$ pasos a $S$.
    \item Como $M$ hace $O(t(n))$ pasos en total, $S$ hace $O(t(n)) \cdot O(t(n)) = O(t(n)^2)$ pasos.
\end{itemize}

\begin{idea}
El factor cuadrático es \emph{polinomial}: si $t(n)$ era un polinomio, $t(n)^2$ también lo es. Por lo tanto, la clase $P$ no cambia al pasar de \MTs\ multicinta a una sola cinta.
\end{idea}

\subsection{Tesis de Church-Turing Extendida}

\begin{tcolorbox}[colback=yellow!10!white, colframe=orange!70!black, title=\textbf{Tesis de Church-Turing Extendida}, fonttitle=\bfseries]
Si un problema es decidible en tiempo $O(t(n))$ por \emph{algún} modelo razonable de computación \textbf{determinista}, entonces es decidible en $O(t(n)^k)$ por \emph{cualquier} otro modelo razonable determinista, para algún $k$.
\end{tcolorbox}

\textbf{Consecuencia:} la clase $P$ es \emph{robusta} en sentido fuerte. No depende de si elegimos máquinas de Turing, máquinas RAM, lambda cálculo, o un lenguaje de programación moderno: el conjunto de problemas resolubles en tiempo polinomial es el mismo.

\begin{observacion}
La tesis tiene la palabra \textbf{determinista} subrayada. El caso \emph{no determinista} es distinto: la simulación de una \MT\ no determinista por una determinista cuesta \emph{exponencial}, no polinomial. Esta brecha entre determinismo y no determinismo es la raíz del famoso problema $P$ vs.\ $NP$ (que se aborda en los teóricos posteriores).
\end{observacion}

\section{Tiempo en máquinas de Turing no deterministas}

\subsection{Definición}

Las máquinas de Turing no deterministas pueden tener varios sucesores en cada configuración. Por eso, la ejecución sobre una entrada forma un \emph{árbol}, no una secuencia lineal.

\begin{definicion}[Tiempo en \MT\ no determinista]
El tiempo de ejecución de una \MT\ no determinista $N$ sobre una entrada de tamaño $n$ es el largo de la rama \emph{más larga} del árbol de cómputo. Es decir, $T_N(n) \leq t(n)$ significa que \emph{todas} las ramas posibles tienen a lo sumo $t(n)$ pasos.
\end{definicion}

\subsection{Simulación: no determinista $\Rightarrow$ determinista}

\begin{teorema}
Si una \MT\ no determinista $N$ decide un problema en $O(t(n))$ pasos, entonces existe una \MT\ determinista $M$ que lo decide en $2^{O(t(n))}$ pasos.
\end{teorema}

\textbf{Idea (BFS sobre el árbol de ejecución).}
\begin{itemize}
    \item Sea $b$ el factor de ramificación máximo (cantidad máxima de sucesores de una configuración). Como $\delta(q, a)$ es un conjunto finito, $b$ está acotado por una constante.
    \item El árbol de ejecución tiene profundidad $\leq t(n)$, así que su cantidad total de nodos es a lo sumo $b^{t(n)}$.
    \item Si $b = 2^c$ para alguna constante $c$, entonces $b^{t(n)} = 2^{c \cdot t(n)} \in 2^{O(t(n))}$.
\end{itemize}

La máquina determinista $M$ recorre todo el árbol con BFS y acepta si y solo si alguna rama acepta. Como debe visitar a lo sumo todos los nodos del árbol, el costo es $2^{O(t(n))}$.

\begin{intuicion}
La simulación \emph{no es polinomial}: pasar de tiempo $O(t(n))$ no determinista a determinista cuesta una explosión exponencial. Este es el motivo por el cual $P$ y $NP$ \emph{podrían} ser clases distintas; la pregunta abierta es si esta simulación es óptima o si existe alguna técnica más inteligente que evite la explosión.
\end{intuicion}

\section{Ejemplos de problemas en $P$}

\subsection{Conectividad en grafos: $\mathit{Path}$}

\begin{ejemplo}
Consideremos
\[
\mathit{Path} = \{ \enc{G, v, w} \mid G \text{ es un grafo con un camino de } v \text{ a } w\}.
\]
\end{ejemplo}

\textbf{Algoritmo ineficiente.} Generar todos los caminos de longitud 1, 2, 3, …, hasta $|V|$, y verificar si alguno conecta $v$ con $w$. La cantidad de posibles caminos es $|V|!$, así que el algoritmo es $O(n!)$. \emph{No tratable.}

\textbf{Algoritmo eficiente (BFS / marcado).}
\begin{enumerate}
    \item Marcar $v$.
    \item Mientras se agreguen nuevos nodos marcados:
    \begin{itemize}
        \item Recorrer las aristas; para cada arista $(a, b)$ con $a$ marcado y $b$ no, marcar $b$.
    \end{itemize}
    \item Si $w$ resultó marcado, aceptar; si no, rechazar.
\end{enumerate}

\textbf{Análisis.} El bucle se ejecuta a lo sumo $|V|$ veces (cada iteración marca al menos un nodo nuevo), y cada pasada recorre las aristas en $O(|E|)$. Total: $O(|V| \cdot |E|) \in O(n^2)$. \textbf{Polinomial}, así que $\mathit{Path} \in P$.

\subsection{Coprimos: aplicación del algoritmo de Euclides}

\begin{ejemplo}
\[
\mathit{Coprimos} = \{ \enc{x, y} \mid \gcd(x, y) = 1\}.
\]
\end{ejemplo}

\textbf{Algoritmo ineficiente (fuerza bruta).} Para cada $z \leq \min(x, y)$, verificar si $z$ divide a ambos. Como los números están codificados en \emph{binario} con $n$ bits, hay $2^n$ valores de $z$ que probar. Tiempo $2^n$: \emph{exponencial} en el tamaño de la entrada. \textbf{No es polinomial.}

\textbf{Algoritmo eficiente (Euclides).}
\begin{enumerate}
    \item Repetir hasta que $y = 0$:
    \begin{itemize}
        \item $x \leftarrow x \bmod y$.
        \item Intercambiar $x$ y $y$.
    \end{itemize}
    \item Si $x = 1$, aceptar (son coprimos); si no, rechazar.
\end{enumerate}

\textbf{Análisis.}
\begin{itemize}
    \item En cada iteración, $x$ se reduce \emph{al menos a la mitad}.
    \item Si los números tienen $n$ bits, alcanzan $O(n)$ iteraciones para terminar.
    \item Cada operación módulo se computa en tiempo polinomial $p(n)$.
\end{itemize}

Total: $O(n \cdot p(n))$, \textbf{polinomial}.

\begin{idea}
Atención al tamaño de la entrada: para problemas numéricos, ``tamaño de la entrada'' significa \emph{cantidad de bits}, no el valor del número. Un algoritmo ``en $x$ pasos'' con $x$ del orden de $2^n$ \emph{no} es polinomial. Por eso la fuerza bruta sobre divisores no califica.
\end{idea}

\subsection{Pertenencia para gramáticas libres de contexto (CYK)}

\begin{ejemplo}
\[
\mathit{GLC} = \{ \enc{G, w} \mid w \text{ es generado por la gramática libre de contexto } G\}.
\]
\end{ejemplo}

\textbf{Idea.} Usar \emph{programación dinámica} y la \textbf{Forma Normal de Chomsky}, donde las producciones son $A \to BC$ o $A \to a$ (o $A \to \varepsilon$ para la cadena vacía).

\textbf{Estructura de datos.} Dado $w = w_1 w_2 \cdots w_n$, construimos una matriz $M[i][j]$, donde $M[i][j]$ es el conjunto de variables que pueden generar la subcadena $w_i \cdots w_j$.

\textbf{Algoritmo (CYK).}
\begin{enumerate}
    \item Si $w = \varepsilon$: aceptar sii $S \to \varepsilon$ es una regla.
    \item Para $i = 1, \dots, n$ y cada variable $A$:
    \begin{itemize}
        \item Si $A \to w_i$ es regla, agregar $A$ a $M[i][i]$.
    \end{itemize}
    \item Para $l = 2, \dots, n$ (longitud de la subcadena):
    \begin{itemize}
        \item Para $i = 1, \dots, n - l + 1$ y $j = i + l - 1$:
        \item Para $k = i, \dots, j - 1$ y cada regla $A \to BC$:
        \begin{itemize}
            \item Si $B \in M[i][k]$ y $C \in M[k+1][j]$, agregar $A$ a $M[i][j]$.
        \end{itemize}
    \end{itemize}
    \item Aceptar si $S \in M[1][n]$.
\end{enumerate}

\textbf{Análisis.}
\begin{itemize}
    \item Los bucles anidados sobre $l, i, k$ y reglas hacen $O(|V| \cdot n^3)$ trabajo, donde $|V|$ es la cantidad de variables. Es \textbf{polinomial}.
\end{itemize}

Por lo tanto $\mathit{GLC} \in P$.

\begin{intuicion}
El algoritmo CYK es un caso bonito de \emph{programación dinámica}: en lugar de explorar todas las derivaciones (lo que sería exponencial), reutilizamos resultados anteriores. La forma normal de Chomsky es la herramienta que permite ``trocear'' las derivaciones en sub-derivaciones más pequeñas.
\end{intuicion}

\section{Resumen final}

Ideas centrales para llevar:

\begin{enumerate}
    \item \textbf{Complejidad vs.\ computabilidad.} Computabilidad pregunta ``¿se puede?''; complejidad pregunta ``¿se puede \emph{rápido}?''. En complejidad ya trabajamos con problemas computables.

    \item \textbf{Tiempo como función.} $T_M : \N \to \R^{\geq 0}$ mide el peor caso. Lo que interesa es el comportamiento asintótico, no las constantes.

    \item \textbf{Notaciones asintóticas.}
    \begin{itemize}
        \item $O(g)$: cota superior. $f \leq c \cdot g$ asintóticamente.
        \item $o(g)$: cota superior estricta. $\lim f/g = 0$.
        \item $\Omega(g)$: cota inferior.
        \item $\Theta(g)$: mismo orden de crecimiento (ambas).
    \end{itemize}

    \item \textbf{Análisis cuidadoso paga.} El mismo problema $\{0^n 1^n\}$ se resuelve en $O(n^2)$ con un algoritmo ingenuo y en $O(n \log n)$ con una idea de ``dividir a la mitad''.

    \item \textbf{Clases de complejidad.} $\mathrm{Time}(t(n))$ agrupa los problemas decidibles en $O(t(n))$ pasos por una \MT\ determinista.

    \item \textbf{La clase $P$.} $P = \bigcup_{k > 0} \mathrm{Time}(O(n^k))$: los problemas resolubles en tiempo polinomial. Es la formalización de ``problemas tratables''. Los polinomios escalan bien; los exponenciales no.

    \item \textbf{Simulación polinomial entre modelos deterministas.} \MTs\ multicinta se simulan en una cinta con costo $O(t(n)^2)$. Esto es polinomial, así que $P$ no depende del modelo (determinista) elegido.

    \item \textbf{Tesis de Church-Turing Extendida.} Todos los modelos razonables \emph{deterministas} se simulan entre sí en tiempo polinomial. Por lo tanto, la noción de ``algoritmo polinomial'' es invariante de modelo.

    \item \textbf{El no determinismo cambia las reglas.} Una \MT\ no determinista de tiempo $t(n)$ se simula con una determinista de tiempo $2^{O(t(n))}$: explosión exponencial. Esta brecha motiva la clase $NP$ y el famoso problema $P \stackrel{?}{=} NP$.

    \item \textbf{Ejemplos de $P$.} $\mathit{Path}$ (BFS, $O(|V| \cdot |E|)$), $\mathit{Coprimos}$ (Euclides, $O(n \cdot p(n))$), $\mathit{GLC}$ (CYK, $O(|V| \cdot n^3)$). \textbf{Importante:} el tamaño de la entrada en problemas numéricos es la cantidad de \emph{bits}, no el valor del número.
\end{enumerate}

\begin{intuicion}
La clase $P$ es el punto de partida de toda la complejidad computacional. Definir ``eficiente'' como ``polinomial'' no es perfecto (un algoritmo $O(n^{100})$ es polinomial pero impráctico), pero captura mucho mejor que cualquier otra alternativa la intuición de que ``los buenos algoritmos escalan''. A partir de $P$ se construye la jerarquía de clases ($NP$, $\mathrm{PSPACE}$, etc.), las nociones de \emph{reducción polinomial} y, eventualmente, las preguntas que aún hoy permanecen abiertas en la teoría de la computación.
\end{intuicion}

\end{document}
